% Vista explodida do invólucro (exemplo fictício), em projeção oblíqua:
% a base com a abraçadeira, a placa e a tampa com as ranhuras.
\begin{tikzpicture}[x={(7.2mm,0mm)}, y={(3.1mm,2.2mm)}, z={(0mm,7.2mm)},
    line join=round, font=\scriptsize,
    aresta/.style={draw=nro-tinta, line width=0.45pt}]
  % uma caixa: canto (x,y,z), largura (x), profundidade (y), altura (z), cores
  \def\caixa#1#2#3#4#5#6#7{%
    \filldraw[aresta, fill=#7!70] (#1+#4,#2,#3) -- (#1+#4,#2+#5,#3) -- (#1+#4,#2+#5,#3+#6) -- (#1+#4,#2,#3+#6) -- cycle;
    \filldraw[aresta, fill=#7!45] (#1,#2,#3) -- (#1+#4,#2,#3) -- (#1+#4,#2,#3+#6) -- (#1,#2,#3+#6) -- cycle;
    \filldraw[aresta, fill=#7!20] (#1,#2,#3+#6) -- (#1+#4,#2,#3+#6) -- (#1+#4,#2+#5,#3+#6) -- (#1,#2+#5,#3+#6) -- cycle;}
  % a abraçadeira, sob a base
  \filldraw[aresta, fill=nro-cinza!30] (4.2,4,-1.3) -- (7.8,4,-1.3) -- (7.8,4,0) -- (4.2,4,0) -- cycle;
  \draw[aresta, fill=white] (6,4,-1.9) circle (1.9mm);
  \node[anchor=west, align=left] at (13.2,4,-1.2) {abraçadeira para tubo de 35\,mm,\\braço de 8\,mm e calço de TPU};
  % a base (12 x 9 x 1,8)
  \caixa{0}{0}{0}{12}{9}{1.8}{nro-fundo}
  \foreach \x/\y in {1/1,11/1,1/8,11/8} \fill[nro-synapse!80!black] (\x,\y,1.8) circle (0.9mm);
  \node[anchor=west, align=left] at (13.2,4.5,0.9) {base em PETG, paredes de 2,4\,mm,\\insertos de latão M3};
  % as linhas de montagem (antes da placa e da tampa, que as escondem)
  \foreach \x/\y in {1/1,11/1,1/8,11/8} \draw[nro-claro, densely dashed] (\x,\y,1.8) -- (\x,\y,7.8);
  % a placa (10 x 7,5 x 0,2), com os componentes altos
  \caixa{1}{0.75}{4.2}{10}{7.5}{0.25}{nro-axon}
  \caixa{2}{2}{4.45}{2.5}{2.5}{0.6}{nro-cinza}
  \caixa{6}{1.2}{4.45}{4}{1.2}{0.7}{nro-tinta}
  \caixa{9.6}{3}{4.45}{1.2}{3}{0.9}{nro-synapse}
  \node[anchor=west, align=left] at (13.2,4.5,4.6) {placa de potência, revisão B\\(Atividade 1)};
  % a tampa, com as ranhuras em labirinto
  \caixa{0}{0}{7.8}{12}{9}{1.6}{nro-fundo}
  \foreach \k in {0,...,5} \filldraw[draw=nro-tinta, line width=0.3pt, fill=nro-tinta!55]
    (2+\k*1.3,5,9.4) -- (2.6+\k*1.3,5,9.4) -- (2.6+\k*1.3,8,9.4) -- (2+\k*1.3,8,9.4) -- cycle;
  \draw[aresta, fill=nro-pulso!80] (3.2,2.2,9.4) circle (2.2mm);
  \node[anchor=west, align=left] at (13.2,4.5,8.7) {tampa em PETG com ranhuras em\\labirinto e o botão de parada};
\end{tikzpicture}
